\documentclass[12pt, a4paper]{article}
\usepackage[utf8]{inputenc}
\usepackage{geometry}
\geometry{a4paper, margin=1in}
\usepackage{hyperref}
\usepackage{graphicx}
\usepackage{listings}
\usepackage{xcolor}
\usepackage{amsmath}
\usepackage{amsfonts}

\definecolor{codegreen}{rgb}{0,0.6,0}
\definecolor{codegray}{rgb}{0.5,0.5,0.5}
\definecolor{codepurple}{rgb}{0.58,0,0.82}
\definecolor{backcolour}{rgb}{0.95,0.95,0.92}

\lstset{
    backgroundcolor=\color{backcolour},   
    commentstyle=\color{codegreen},
    keywordstyle=\color{magenta},
    numberstyle=\tiny\color{codegray},
    stringstyle=\color{codepurple},
    basicstyle=\ttfamily\footnotesize,
    breakatwhitespace=false,         
    breaklines=true,                 
    captionpos=b,                    
    keepspaces=true,                 
    numbers=left,                    
    numbersep=5pt,                  
    showspaces=false,                
    showstringspaces=false,
    showtabs=false,                  
    tabsize=2
}

\title{\textbf{Advanced Autonomous Mapping and Navigation Simulation for the QUANSER QBot 2 for QUARC}}
\author{Research and Implementation Documentation}
\date{\today}

\begin{document}

\maketitle

\begin{abstract}
This document details the development of a high-fidelity 3D simulation and research platform for the QUANSER QBot 2 differential-drive robot. The project integrates real-time SLAM-inspired mapping, persistent point cloud accumulation, and A* pathfinding algorithms within a modern web-based dashboard architecture. The system enables robust testing of autonomous navigation logic in complex maze environments prior to hardware deployment.
\end{abstract}

\tableofcontents
\newpage

\section{Introduction}
The transition from quadrupedal platforms to differential-drive wheeled robots, such as the QUANSER QBot 2, presents unique challenges in odometry and trajectory planning. This research focuses on building a unified simulation environment that bridges the gap between high-level Python control logic and low-level 3D visualization.

\section{System Architecture}
The system architecture follows a decoupled model-view-controller (MVC) pattern, utilizing WebSockets for low-latency bidirectional telemetry.

\subsection{Differential Drive Kinematics}
The QBot 2 is modeled as a differential drive robot. Its motion is defined by the independent velocities of the left ($v_L$) and right ($v_R$) wheels. The linear velocity $v$ and angular velocity $\omega$ of the robot's center of mass are:
\begin{equation}
    v = \frac{v_R + v_L}{2}, \quad \omega = \frac{v_R - v_L}{L}
\end{equation}
Where $L$ is the track width (distance between wheels). The state $[x, y, \theta]$ is updated using:
\begin{equation}
    x_{t+1} = x_t + v \cos(\theta) \Delta t, \quad y_{t+1} = y_t + v \sin(\theta) \Delta t, \quad \theta_{t+1} = \theta_t + \omega \Delta t
\end{equation}
This kinematic model ensures that the simulation accurately reflects the movement constraints of the physical QBot 2 hardware.

\subsection{Backend: FastAPI and WebSocket Relay}
The backend serves as a high-speed relay between the Python control script and the 3D visualization engine. Using \texttt{FastAPI} and \texttt{Uvicorn}, the server handles asynchronous WebSocket connections (\texttt{ws://0.0.0.0:8000/api/ws/sim}), allowing the robot to transmit its state and receive movement vectors in real-time.

\subsection{Frontend: Three.js Visualization Engine}
The frontend leverages \texttt{Three.js} for hardware-accelerated 3D rendering. It consists of three primary rendering targets:
\begin{enumerate}
    \item \textbf{Main Perspective View}: An overview of the maze and robot.
    \item \textbf{FPV Camera}: A First-Person View rendered from the QBot 2's sensor perspective.
    \item \textbf{3D Mapping View}: A dedicated secondary scene for SLAM reconstruction.
\end{enumerate}

\section{3D Mapping and SLAM Reconstruction}
The core of the robot's spatial awareness is its mapping capability, which is divided into two distinct methodologies: Point Cloud accumulation and Voxel reconstruction.

\subsection{LiDAR Raycasting Emulation}
Spatial data is acquired through virtual LiDAR emulation. The system performs 144 radial raycasts (2.5° resolution) every frame. The intersection point $\mathbf{P}$ of each ray with the environment is calculated as:
\begin{equation}
    \mathbf{P} = \mathbf{O} + d \cdot \vec{\mathbf{D}}
\end{equation}
Where $\mathbf{O}$ is the sensor origin, $\vec{\mathbf{D}}$ is the normalized direction vector, and $d$ is the distance to the nearest collision.

\subsection{Persistent Point Cloud}
Each detected point $\mathbf{P}$ is added to a persistent buffer of $100,000$ points. This allows the robot to "remember" the structure of the room even after it has moved past a specific area. This data is rendered as a cyan point cloud in the dedicated Mapping Tab.

\subsection{Voxel Environment Reconstruction}
To provide a solid structural map, the system implements voxelization. The 3D space is divided into a grid of size $0.5m^3$. When a LiDAR point is detected within a grid cell $(v_x, v_y, v_z)$, the cell is marked as occupied:
\begin{equation}
    v_x = \lfloor x / s \rfloor, \quad v_y = \lfloor y / s \rfloor, \quad v_z = \lfloor z / s \rfloor
\end{equation}
\subsection{Simultaneous Localization and Mapping (SLAM)}
The platform implements a simplified SLAM approach where localization is provided by simulated high-precision odometry, while mapping is performed through LiDAR integration.
\begin{itemize}
    \item \textbf{Global Consistency:} Points are transformed from the robot's local frame $\mathbf{L}$ to the world frame $\mathbf{W}$ using the transformation matrix:
    \begin{equation}
        \mathbf{P}_W = \mathbf{T}_{R \to W} \mathbf{P}_L
    \end{equation}
    \item \textbf{Noise Filtering:} The voxelization process acts as a spatial low-pass filter, reducing sensor noise and memory overhead by quantizing points into a discrete occupancy grid.
\end{itemize}

Where $s$ is the voxel size. Occupied cells are rendered using an \texttt{InstancedMesh} to ensure high performance while maintaining a solid visual representation of the mapped world.

\section{Navigation and Path Planning}
The QBot 2 utilizes the A* (A-Star) algorithm for collision-free routing through the maze.

\subsection{A* Algorithm Implementation}
The pathfinder calculates the optimal path from current coordinates to a waypoint by minimizing the cost function:
\begin{equation}
    f(n) = g(n) + h(n)
\end{equation}
Where:
\begin{itemize}
    \item $g(n)$ is the cost of the path from the start node to node $n$.
    \item $h(n)$ is a heuristic estimate of the cost from $n$ to the goal (Manhattan distance).
\end{itemize}

\subsection{Search Space Discretization}
The maze is discretized into a $20 \times 20$ navigation grid. Each cell is evaluated for its "traversability" based on the presence of wall meshes.
\begin{itemize}
    \item \textbf{Open Set:} A priority queue of nodes to be evaluated, sorted by $f(n)$.
    \item \textbf{Closed Set:} Nodes already visited to prevent redundant cycles and infinite loops.
    \item \textbf{Path Smoothing:} Once a grid-based path is found, the robot performs linear interpolation between nodes for smooth transitions.
\end{itemize}

\subsection{Collision Avoidance}
Beyond global pathfinding, the robot maintains a local reactive layer. If the forward LiDAR distance drops below a safety threshold ($0.3m$), an emergency stop is triggered, and the global path is recalculated to avoid the obstruction.

\section{Control API: qbot\_sim.py}
A standalone Python SDK was developed to allow researchers to control the robot with minimal code. The API handles the complexity of asynchronous communication and telemetry parsing.

\begin{lstlisting}[language=Python, caption=Sample Control Script]
from qbot_sim import QBot2Sim

robot = QBot2Sim()
robot.reset()
robot.log("Starting Patrol...")

# Navigate to a specific coordinate
robot.move_to(17.0, -17.0) 

# Access sensor data
pos = robot.get_position()
lidar = robot.get_lidar_front()
\end{lstlisting}

\section{Real-World Applications}
The simulation platform is designed for several high-impact research areas:
\begin{itemize}
    \item \textbf{Autonomous Surveillance}: Training robots to detect anomalies in industrial corridors.
    \item \textbf{Hazardous Environment Mapping}: Building maps in areas inaccessible or dangerous for humans.
    \item \textbf{Algorithmic Validation}: Rapid prototyping of new SLAM and pathfinding algorithms without the risk of hardware damage.
\end{itemize}

\section{Conclusion}
The QUANSER QBot 2 simulation platform provides a robust environment for autonomous systems research. By separating the 3D mapping reconstruction into a dedicated SLAM view and implementing advanced pathfinding, the platform serves as a production-grade validation tool for complex robotic behaviors.

\end{document}
